\begin{enumerate}

\item La inercia de la sección idealizada es $4 B y^2$.

\item El flujo en cada alma es el salto en el larguero superior.

\item La tensión tangencial es el flujo dividido entre el espesor.

\end{enumerate}

\noindent\textbf{Operaciones y resultados intermedios.}

\begin{enumerate}

\item Inercia.

\[\fitmath{y_{1}=4\cdot 400\cdot {100}^{2}\approx 1{,}6\times10^{7}\,\mathrm{mm^{4}}}\]

\item Flujo en las almas.

\[\fitmath{y_{2}=\frac{10000\cdot 400\cdot 100}{4\cdot 400\cdot {100}^{2}}\approx 25\,\mathrm{N/mm}}\]

\item Tensión tangencial.

\[\fitmath{y_{3}=\frac{\frac{10000\cdot 400\cdot 100}{4\cdot 400\cdot {100}^{2}}}{2}\approx 12{,}5\,\mathrm{MPa}}\]

\end{enumerate}
